\section{Proof of the volume theorem}
\label{sec:conclusion}

\begin{proof}[Proof of Theorem~\ref{thm:volume}]
We induct on \(m\ge2\). Fix \(m\), assuming the theorem in dimensions
\(2,\ldots,m-1\); this hypothesis is empty when \(m=2\).
By Proposition~\ref{prop:immersion-reduction}, every counterexample in
dimension \(m\) would be a smooth closed connected embedded minimal
hypersurface.

Suppose such counterexamples exist. Proposition~\ref{prop:attainment}
gives one of least area \(H\), denoted by \(\Sigma\), with
\[
\sigma_m<\gamma_m\le H<a_m.
\]
Set \(p=m+3\). Proposition~\ref{prop:index-rigidity} gives
\(\ind(\Sigma)\ge p\). Equality would make \(\Sigma\) a minimal
Clifford product and force \(H\ge a_m\). Therefore
\(\ind(\Sigma)>p\).

By Proposition~\ref{prop:lorentz-family}, there is a detecting family
without mass concentration whose supremum mass \(D_1\) satisfies
\(D_1<H\). Theorem~\ref{thm:minmax} yields a stationary integral
\(m\)-varifold \(V\) on
\(S^{m+1}\) with
\[
\gamma_m\le\|V\|(S^{m+1})<H<a_m.
\]
When \(m=2,3\), choose the smooth embedded min--max output provided by
that theorem, with its integer component multiplicities. Its
Euclidean cone is cyclic modulo two by
Lemma~\ref{lem:cone-operation}. In every higher dimension the cone is
stationary and integral, which are the required hypotheses.

Corollary~\ref{cor:cones-spherical-output}, using the outer induction
hypothesis, now shows that \(V\) is a smooth connected embedded minimal
hypersurface of multiplicity one. Its mass exceeds \(\sigma_m\), so it
is non-equatorial. It is therefore a counterexample with area strictly
smaller than \(H\), a contradiction.

This proves the embedded assertion in dimension \(m\), and
Proposition~\ref{prop:immersion-reduction} gives the full immersion
assertion. The induction proves the theorem for every \(m\ge2\).
\end{proof}
